%% kapa: exam=1 round=30 year=2019 period=1 form=A subject=real_estate subject_ko=부동산학원론 seq=37 number=117
%% source: Q-Net 감정평가사 기출문제 (cst00309) artlSeq=5205205 file=extracted/1008/5205205/2019년 제30회 감정평가사 1차시험 문제지/2019년 제30회 감정평가사 1차시험 문제지/2019년도 제30회 감정평가사 1차 1교시 A형.hwp
%% conversion: kordoc 4.24.51 -> tools/gen1.py (auto); review=proofread-sonnet; flags=none
\begin{question}{117}
부동산 중개계약에 관한 설명으로 옳은 것은?
\begin{choices}
  \item 순가중개계약은 중개의뢰인이 다수의 개업공인중개사에게 의뢰하는 계약의 형태이다.
  \item 독점중개계약을 체결한 개업공인중개사는 자신이 거래를 성립시키지 않았을 경우 중개보수를 받지 못한다.
  \item 전속중개계약을 체결한 개업공인중개사는 누가 거래를 성립시켰는지에 상관없이 중개보수를 받을 수 있다.
  \item 공동중개계약은 다수의 개업공인중개사가 상호 협동하여 공동으로 중개 역할을 하는 것이다.
  \item 일반중개계약은 거래가격을 정하여 개업공인중개사에게 제시하고, 이를 초과한 가격으로 거래가 이루어진 경우 그 초과액을 개업공인중개사가 중개보수로 획득하는 방법이다.
\end{choices}
\end{question}
\answer{4}
